\section{Eliminating the introduction witness}
\label{sec:logical-proof}

Write $\mathcal N$ for the pinned eight-rule native calculus and
$\mathcal K$ for the calculus in which case introduction always chooses
the newest numerical parameter, or $0$ in an empty numerical scope.
The two systems have exactly the same terms, formulas, contexts, and
ordinary readings. In particular, $\mathcal K$ retains the separate
progress rule for literal equations. The result in this section does not
require merging the two introduction modes.

\begin{theorem}[Exact equivalence on open sequents]
\label{thm:newest-equivalence}
For every numerical scope, every finite native context $\Gamma$, and
every native formula $A$,
\[
 \Gamma\vdash_{\mathcal K} A
 \quad\Longleftrightarrow\quad
 \Gamma\vdash_{\mathcal N} A.
\]
The forward translation preserves every proof instruction except that
it supplies the witness selected by $\mathcal K$. The reverse translation
is effective and preserves every term and formula of the end-sequent.
\end{theorem}

\subsection{Why a proof must be rebuilt under a new parameter}

The selected witness depends on the numerical scope, including parameters
that do not occur in the conclusion. Consequently, adding a fresh parameter
to an existing proof tree can change the witness required by an introduction
instruction. Numerical weakening is therefore not an immediate operation on
$\mathcal K$ proof trees. The usual proof of cut also cannot be assumed:
when the receiving proof enters a fresh-parameter branch, the proof being
inserted would have to be weakened numerically.

We use \emph{persistent proof builders}. Such a builder for
$\Gamma\vdash A$ supplies, for every extension of the numerical scope,
a newly constructed $\mathcal K$ proof of the shifted $A$ from the shifted
$\Gamma$. A builder may inspect the extended scope and generate its own
fresh parameters. It does not shift a previously constructed proof tree.
All builders below are uniform in the displayed term and formula arguments.

Context weakening at a fixed numerical scope is admissible by induction on
proofs: adding formula assumptions does not change the selected numerical
witness, either at the root or inside a branch. This gives the following
limited cut principle before general cut has been established.

\begin{lemma}[Persistent cut]
\label{lem:persistent-cut}
A persistent builder of $\Gamma\vdash A$ may replace an assumption $A$
in any $\mathcal K$ derivation from $\Gamma,A$. If that second derivation
is also supplied persistently, the result is persistent.
\end{lemma}

\begin{proof}
Traverse the second proof tree. At an occurrence of the replaced
assumption below $k$ fresh numerical parameters, call the first builder
at extension $k$ and context-weaken its result by the local assumptions.
Every retained instruction has exactly its former numerical scope, so
every selected witness is unchanged. For a persistent receiving proof,
perform this traversal anew at each requested ambient extension.
\end{proof}

We will use this lemma only with the explicitly uniform builders
constructed below. It does not assert that every $\mathcal K$ derivation
is already persistent.

\subsection{Closure introduction and formula transport}

The selector term is
\[
 C(v,a,b)=I(v,a;u.b),\qquad
 C(0,a,b)=a,\qquad C(Sv,a,b)=b.
\]
The unused update variable in this definition is intentional.

\begin{lemma}[Uniform closure introduction]
\label{lem:newest-qcase}
The following rule is admissible, with $z$ fresh:
\[
 \frac{\Gamma,B(0)\vdash_{\mathcal K} B(Sz)}
      {\Gamma\vdash_{\mathcal K}Q_yB(y)}.
\]
Its construction works in every numerical scope. Hence a base proof
$\Gamma\vdash B(0)$ and a successor proof $\Gamma\vdash B(Sz)$ give
$\Gamma\vdash B(t)$ for any term $t$.
\end{lemma}

\begin{proof}
Introduce $D_{x,y}B(y)$, where $x$ has no occurrence. Substituting the
automatically selected witness for $x$ therefore leaves the premise
unchanged. Eliminate this $D$ formula: its opened assumption is precisely
$Q_yB(y)$, independent of the fresh witness, and that assumption may be
exported. For the last assertion, context-weaken the successor proof by
$B(0)$, apply the first assertion, and use QE at $t$.

When this construction is used inside a builder, its successor proof is
generated with the $z$ fresh in that invocation. There is no assertion
that a fixed successor proof can be arbitrarily substituted or weakened.
\end{proof}

\begin{lemma}[Uniform Leibniz rule]
\label{lem:newest-transport}
For every native formula template $F(v)$ and terms $s,t$, there is a
persistent builder of
\[
 s=t,\ F(s)\ \vdash_{\mathcal K}\ F(t).
\]
Consequently persistent proofs of $s=t$ and $F(s)$ give a persistent
proof of $F(t)$ over the same context.
\end{lemma}

\begin{proof}
Induct on the formula template. For an equation, the crossed one-term
transport rule gives the ordinary equality operations. Transport in
the left and right term templates takes two applications. These are
uniform finite proof recipes.

For $F(v)=Q_yB(v,y)$, retain $s=t$ and $Q_yB(s,y)$ as assumptions.
To introduce the target closure, assume $B(t,0)$ in its fresh-$z$
premise. The induction hypothesis, with the equality reversed,
gives $B(s,0)$. QE applied to the retained source closure gives
$B(s,Sz)$, and the induction hypothesis gives $B(t,Sz)$.
Lemma~\ref{lem:newest-qcase} finishes the construction. Each matrix
transport has smaller formula complexity; its uniform builder can be
used at the current terms and scope. Compositions use
Lemma~\ref{lem:persistent-cut}.

For $F(v)=D_{u,y}B(v,u,y)$, eliminate the retained source formula at a
fresh parameter $b$. Its opened assumption is $Q_yB(s,b,y)$.
Introduce the target $D_{u,y}B(t,u,y)$ in this context. The parameter
$b$ is now newest, so it is exactly the witness required by the rule.
Under $B(t,b,0)$, the matrix induction hypothesis transports to
$B(s,b,0)$; QE on the opened source closure gives $B(s,b,Sz)$;
the matrix induction hypothesis transports back to $B(t,b,Sz)$.
The target $D$ formula does not contain $b$, so DE discharges it.

In every invocation after an ambient scope extension, a new $b$ is
generated \emph{after} all ambient parameters. Thus the same recipe
still introduces at the newest parameter. This proves persistence
without numerically weakening a proof containing an old $b$.
The final assertion follows by persistent cut.
\end{proof}

\subsection{The arithmetic dependency}

The only arithmetic fact needed for naming is a scalar predecessor term
$P$ with a uniform proof of
\begin{equation}
 P(0)=0,\qquad P(Sr)=r.
 \label{eq:newest-predecessor}
\end{equation}
The scalar bootstrap supplied with the comparison material constructs
such a finite term using $0,S,I$ alone. Its resources are computation,
equality, induction on a literal equation, and numerical cases with
equation conclusions. The accompanying arithmetic verification explains
how to carry that bootstrap out before the witness-elimination theorem.
In particular it does not appeal to general existential introduction.

For completeness, even conditional equation cases can be implemented
uniformly at this stage. For equations $E\equiv p=q$ and $F\equiv r=s$,
define
\[
 E\Rightarrow_{\rm eq} F
 :=Q_v\bigl(C(v,p,r)=C(v,q,s)\bigr).
\]
Its introduction rule follows from Lemma~\ref{lem:newest-qcase}, because
the zero row computes to $E$ and every successor row computes to $F$.
Its elimination rule is QE at $1=S0$, followed by computation.
Now suppose a numerical case argument derives $F$ under $e=0$, and
under $e=Sz$ with $z$ fresh. Apply these equation-implication rules to
\[
 H(y):=(e=y)\Rightarrow_{\rm eq}F.
\]
The two branches give $H(0)$ and $H(Sz)$. Clean numerical cases give
$H(e)$, which applied to reflexivity $e=e$ gives $F$. With branch
builders this is a persistent construction, including arbitrary retained
side assumptions. Thus contextual equation cases in the written scalar
bootstrap create no dependence on the theorem being proved.

Applying $P$ to an equation $Ss=St$, and using
\eqref{eq:newest-predecessor}, gives the persistent rule
\begin{equation}
 Ss=St\ \vdash_{\mathcal K}\ s=t.
 \label{eq:newest-injective}
\end{equation}
No naming packet is used to establish this rule.

\subsection{A fresh name for every term}

\begin{lemma}[Naming without choosing a witness]
\label{lem:newest-name}
Let
\begin{equation}
 X(y):=D_{x,v}
 \bigl(C(v,0,C(y,0,Sx))=C(v,0,y)\bigr).
 \label{eq:newest-packet}
\end{equation}
For every term $t$, $X(St)$ has a persistent proof. When this formula
is opened at a fresh parameter $b$, its retained closure yields a
persistent proof of $b=t$.
\end{lemma}

\begin{proof}
First prove $X(0)$. Whatever witness $w$ the current scope selects,
the case-introduction successor premise computes to
$C(0,0,Sw)=0$, hence to $0=0$. The assumption in that premise is
unneeded. Therefore $X(0)$ is available even in an empty scope, and
also when the selected witness is an unrelated ambient parameter.

For $X(Sz)$, generate the numerical-case parameter $z$ freshly.
It is the newest parameter when the inner $D$ is introduced, so the
rule selects $x=z$. The inner successor premise computes to
$C(Sz,0,Sz)=Sz$, hence to $Sz=Sz$. Again the premise's base assumption
is unneeded. Lemma~\ref{lem:newest-qcase} now gives $X(r)$ for an
arbitrary target term $r$; take $r=St$.

Open this packet at a fresh $b$. Its closure has body
\[
 C(v,0,C(St,0,Sb))=C(v,0,St).
\]
At $v=0$ the body is a computational reflexivity. QE at $1$, followed
by selector computation, gives $Sb=St$. Equation
\eqref{eq:newest-injective} gives $b=t$.

All these proofs are rebuilt uniformly at each ambient scope. In
particular the $z$ used to prove the successor case is always generated
after the current parameters. The last alias proof uses only the retained
closure, QE, computation, equality and the predecessor builder, so it
remains persistent under further parameters.
\end{proof}

The ordinary reading of $X(y)$ is
$\exists x\,[C(y,0,Sx)=y]$. The argument uses this formula at $St$,
where the selector computes away, and then uses injectivity. It never
attempts to introduce $\exists x\,(x=t)$ directly with the prohibited
witness $t$.

\subsection{The protected compiler}

\begin{proof}[Proof of Theorem~\ref{thm:newest-equivalence}]
The forward direction is immediate: every restricted introduction is an
instance of the old introduction rule. For the converse, strengthen the
induction statement on an old derivation $p$ of $\Gamma\vdash A$.
The compiler accepts:

\begin{enumerate}
 \item an arbitrary term substitution $\sigma$ from the source numerical
       scope into an arbitrary target scope;
 \item a target context $\Delta$;
 \item for each $E\in\Gamma$, a persistent $\Delta$-proof builder of
       $E[\sigma]$.
\end{enumerate}

It returns a persistent $\Delta$-proof builder of $A[\sigma]$.
To answer a request at an extended target scope, it reruns its construction
with shifted $\sigma$ and the supplied hypothesis builders. It does not
shift its previously returned proof.

Hypothesis, computation, one-term transport and QE are handled directly.
For a source QI-progress instruction, generate its fresh target parameter,
extend $\sigma$ by that parameter, and invoke the induction hypothesis
on the step subderivation. Supply the old hypothesis builders in the
extended scope and the persistent literal-hypothesis builder for its
equation assumption. Apply QI-progress. Source DE is treated in the
same manner: open its translated major premise at a fresh parameter,
extend $\sigma$, and translate the minor subderivation using the opened
closure as a new persistent hypothesis. Every eigenvariable condition
is preserved.

Consider source case introduction with arbitrary witness $t$ and matrix
$B(x,y)$. Put $t'=t[\sigma]$ and let $B'$ be the matrix obtained by
substituting for its ambient free parameters. Build $X(St')$ by
Lemma~\ref{lem:newest-name} and open it at fresh $b$. From the retained
closure obtain the persistent alias $b=t'$.

We now introduce $D_{x,y}B'(x,y)$. At this point $b$ is the newest
parameter, so the target rule selects it. In the rule's fresh-$z$
premise the local hypothesis is $B'(b,0)$. Formula transport gives
a persistent proof of $B'(t',0)$. Invoke the induction hypothesis on
the source case-premise subderivation, using $\sigma$ extended by $z$.
Its original hypotheses are represented by the inherited builders;
its source base hypothesis is represented by this transport builder.
It returns $B'(t',Sz)$. Transport back along $b=t'$ to obtain
$B'(b,Sz)$, apply the restricted introduction, and discharge $b$ using
DE on the naming packet. The exported formula does not contain $b$.

This construction is well founded. Every call to the proof compiler
uses a strict subderivation of $p$. Formula transport uses induction
on formula complexity. The arithmetic and naming builders were
constructed before this compiler, without invoking its equivalence
theorem. Persistent hypotheses, rather than unproved numerical
weakening or general cut, provide all premises needed under new scopes.

Finally use the identity substitution and literal-hypothesis builders.
The output is a $\mathcal K$ proof of the original open sequent.
\end{proof}

\begin{corollary}[Structural rules and the ordinary theory]
The new calculus admits numerical substitution, numerical weakening and
cut. Its ordinary realization is equiconsistent and bi-interpretable
with iPRA, with both numerical domains whole and both numerical
round-trip graphs equal to $x=y$.
\end{corollary}

\begin{proof}
For the structural rules, translate to $\mathcal N$, use its established
structural rule, and translate back. These rules are consequences of
the equivalence, not premises of its proof.

Between $\mathcal K$ and $\mathcal N$ the interpretations are literally
the identity on terms and native formulas. Compose this identity
equivalence with the pinned native-to-ordinary and ordinary-to-iPRA
interpretations. The ordinary meanings of $Q$ and $D$, and the definitions
of the ordinary intuitionistic connectives, are unchanged. Their existing
formula round trips therefore apply unchanged. The numerical maps also
remain unchanged, so their composites retain the identity graph $x=y$.
The designated contradiction is still $0=S0$; its correspondence with
ordinary falsity and its elimination rule are preserved by exact
open-sequent equivalence. Thus consistency is preserved in both
directions as well.
\end{proof}

\paragraph{Verification boundary.}
The argument above is a written metatheorem. The pinned development
provides machine-checked results for the old native calculus and its
iPRA interpretation; those results are not a machine-checked proof of
this new compiler. Executable certificate checks validate their actual
finite certificates only. A complete formalization of the new result
must encode the persistent builders or an equivalent protected
translation, rather than assume numerical weakening of restricted
proof trees.
